\chapter*{Mappa rapida degli esercizi}
\addcontentsline{toc}{chapter}{Mappa rapida degli esercizi}

\begin{longtable}{@{}p{.29\textwidth}p{.29\textwidth}p{.34\textwidth}@{}}
\toprule
\textbf{La traccia chiede di...} & \textbf{Primo metodo da considerare} & \textbf{Capitolo / controllo decisivo}\\
\midrule
\endhead
Determinare un dominio & imporre simultaneamente tutte le condizioni di esistenza & 2--4; denominatori, radici pari, logaritmi, inverse trigonometriche\\
Risolvere un'equazione o disequazione & fissare il dominio, usare trasformazioni equivalenti, verificare i candidati & 2; tabella dei segni e controllo delle soluzioni estranee\\
Studiare o trasformare un grafico & partire dal grafico elementare e applicare traslazioni, scale e riflessioni nell'ordine corretto & 3\\
Risolvere un problema trigonometrico & ridurre l'angolo, scegliere identità o formula risolutiva, scrivere tutte le soluzioni & 4\\
Operare con numeri complessi & forma algebrica per somme; forma esponenziale per prodotti, potenze e radici & 5\\
Calcolare il limite di una successione & riconoscere dominante, monotonia, ricorrenza, Stolz o sottosuccessioni & 6\\
Calcolare un limite di funzione & sostituzione, algebra, equivalenze, Taylor o de l'Hôpital secondo la forma & 7, 9, 11, 12\\
Verificare continuità o trovare parametri & uguagliare limite e valore; controllare separatamente i laterali & 8\\
Calcolare una derivata o una tangente & dominio, regole di derivazione, catena, valutazione nel punto & 10\\
Applicare un teorema differenziale & verificare tutte le ipotesi prima di usare la conclusione & 11\\
Approssimare una funzione o risolvere una cancellazione & scegliere centro e ordine di Taylor, stimare il resto & 12\\
Eseguire uno studio completo di funzione & seguire la lista di controllo del capitolo, senza saltare dominio e limiti & 13\\
Calcolare una primitiva & riconoscere struttura: immediata, sostituzione, parti, fratti semplici, trigonometrica & 14\\
Calcolare area, valore medio o integrale definito & semplificare con simmetrie; poi primitiva, sostituzione o parti & 15\\
Decidere un integrale improprio & spezzare in ogni singolarità e confrontare localmente con integrali campione & 16\\
Decidere una serie numerica & termine generale, segno, struttura; poi confronto, rapporto/radice, integrale o Leibniz & 17\\
Trovare raggio e intervallo di convergenza & rapporto/radice sui coefficienti e studio separato degli estremi & 18\\
Riconoscere integrali Gamma/Beta & ricondurre mediante sostituzione alle forme standard e usare le ricorrenze & 19\\
\bottomrule
\end{longtable}

